\documentclass[11pt,a4paper]{article}

% ---------- Packages ----------
\usepackage[utf8]{inputenc}
\usepackage[T1]{fontenc}
\usepackage{lmodern}
\usepackage[a4paper,margin=0.82in]{geometry}
\usepackage{microtype}
\usepackage{graphicx}
\usepackage{xcolor}
\usepackage{tabularx}
\usepackage{array}
\usepackage{enumitem}
\usepackage{titlesec}
\usepackage{fancyhdr}
\usepackage{hyperref}
\usepackage{booktabs}
\usepackage{ragged2e}

% ---------- Theme ----------
\definecolor{VVVNavy}{HTML}{101C2C}
\definecolor{VVVCyan}{HTML}{00B8C8}
\definecolor{VVVLight}{HTML}{EAF7F8}
\definecolor{VVVGray}{HTML}{536171}

\hypersetup{
    colorlinks=true,
    linkcolor=VVVNavy,
    urlcolor=VVVCyan,
    pdftitle={Veni Vidi Vici - Event Details},
    pdfauthor={Coding Club, IIT Dharwad}
}

\setlength{\parindent}{0pt}
\setlength{\parskip}{5pt}
\setlist[itemize]{leftmargin=1.5em,itemsep=2pt,topsep=3pt}
\renewcommand{\arraystretch}{1.25}

% ---------- Headings ----------
\titleformat{\section}
  {\Large\bfseries\color{VVVNavy}}
  {\thesection.}{0.55em}{}
  [\vspace{2pt}\color{VVVCyan}\titlerule]

\titleformat{\subsection}
  {\large\bfseries\color{VVVNavy}}
  {}{0pt}{}

% ---------- Header and footer ----------
\pagestyle{fancy}
\fancyhf{}
\lhead{\small\textcolor{VVVGray}{PARSEC 7.0 -- Indian Institute of Technology Dharwad}}
\rhead{\small\textcolor{VVVGray}{VENI VIDI VICI}}
\cfoot{\small\textcolor{VVVGray}{\thepage}}
\renewcommand{\headrulewidth}{0.4pt}
\renewcommand{\headrule}{\hbox to\headwidth{\color{VVVCyan}\leaders\hrule height \headrulewidth\hfill}}

% ---------- Helpers ----------
\newcommand{\domain}[2]{%
    \subsection*{#1}
    #2
}
\newcommand{\placeholder}[1]{\textcolor{red}{[#1]}}

\begin{document}

% ===================== COVER / TITLE =====================
\begin{center}
    {\large\bfseries\color{VVVGray} PARSEC 7.0 -- ANNUAL TECHNICAL FESTIVAL}\\[3pt]
    {\large\bfseries\color{VVVGray} INDIAN INSTITUTE OF TECHNOLOGY DHARWAD}\\[22pt]

    {\Huge\bfseries\color{VVVNavy} VENI VIDI VICI}\\[7pt]
    {\Large\bfseries\color{VVVCyan} Think Like an Attacker. Reason Like a Defender.}\\[16pt]

    \colorbox{VVVLight}{%
        \parbox{0.88\textwidth}{%
            \centering
            \vspace{7pt}
            \textbf{CYBERSECURITY CAPTURE THE FLAG COMPETITION}\\
            Organized by: \textbf{Coding Club, IIT Dharwad}
            \vspace{7pt}
        }%
    }
\end{center}

\vspace{10pt}

\section{Introduction}

PARSEC, the annual technical festival of the Indian Institute of Technology Dharwad,
brings together students, innovators, and technology enthusiasts to explore emerging
technologies and demonstrate their technical abilities.

As part of PARSEC 7.0, \textbf{The Systems Wing} of the Coding Club, Indian Institute of Technology Dharwad, proudly presents
\textbf{Veni Vidi Vici}, a cybersecurity Capture The Flag (CTF) competition designed
to challenge participants across multiple domains of information security.

\textbf{Veni Vidi Vici} aims to go beyond conventional cybersecurity competitions by
introducing participants to \textbf{realistic attack scenarios}, \textbf{carefully designed vulnerable
environments}, and \textbf{multi-stage challenges} that demand analytical thinking, technical
proficiency, and creative problem-solving.

The competition follows a scenario-driven approach, \textbf{encouraging participants to
think like attackers while understanding the underlying principles of defensive
security.} From network investigations and application vulnerabilities to cloud
misconfigurations, cryptographic weaknesses, and emerging AI security threats,
participants will encounter challenges inspired by real-world cybersecurity incidents.

Through an immersive and hands-on experience, \textbf{Veni Vidi Vici} aspires to foster
practical cybersecurity skills, encourage security research, and create an environment
where students can explore the evolving landscape of digital security.

\section{About the Organizers}

The Systems Wing of the Coding Club of IIT Dharwad is a student-driven technical community dedicated
to promoting systems programming, software development, hardware integrations, embedded systems
and emerging areas of computer
science through collaborative learning, technical projects, workshops, and competitive
events.

The wing actively encourages students to explore domains such as:
\begin{itemize}
    \item Systems programming and low-level computing.
    \item Artificial intelligence and machine learning.
    \item Cybersecurity and information security.
    \item Hardware and embedded systems based integrations.
    \item Distributed systems, cloud computing, and networking.
\end{itemize}

Through Veni Vidi Vici, the club seeks to extend its technical initiatives into practical
cybersecurity education by providing participants with an opportunity to engage with
realistic security problems in a controlled environment.

\section{Proposed Event Structure}

\subsection*{Veni Vidi Vici: Capture The Flag Competition}

The competition will feature a collection of carefully curated cybersecurity challenges
distributed across six major domains.

The proposed challenge architecture consists of \textbf{12 primary challenge
environments}, with approximately \textbf{24--30 flags} distributed across different
stages of investigation and exploitation.

Challenges will be organized across three progressive difficulty levels:
\begin{itemize}
    \item \textbf{Foundation:} Fundamental cybersecurity concepts and introductory
    problem-solving.
    \item \textbf{Intermediate:} Multi-step investigations requiring practical technical
    knowledge.
    \item \textbf{Advanced:} Complex attack scenarios requiring deeper analysis,
    exploitation, and lateral reasoning.
\end{itemize}

\subsection*{Challenge Domains}

\domain{1. Network Security}{
Challenges involving packet analysis, network reconnaissance, traffic investigation,
protocol vulnerabilities, and digital forensics.
}

\domain{2. Application Security}{
Exploration of common web application vulnerabilities, authentication weaknesses,
insecure APIs, injection attacks, and access-control failures.
}

\domain{3. Cloud Security}{
Challenges involving cloud infrastructure misconfigurations, identity and access
management, exposed resources, and security weaknesses in cloud environments.
}

\domain{4. System Security and Privilege Escalation}{
Investigation of vulnerable systems, permission misconfigurations, privilege escalation
techniques, and operating-system security.
}

\domain{5. Applied Cryptography}{
Problems involving cryptographic implementation weaknesses, encryption and encoding
mechanisms, cryptanalysis, and secure communication.
}

\domain{6. AI Systems Security}{
Exploration of emerging vulnerabilities in AI-powered systems, including prompt
injection, insecure agent interactions, and weaknesses in AI-integrated applications.
}

\subsection*{Challenge Design Philosophy}

The competition emphasizes realistic, scenario-driven challenges rather than isolated
technical exercises.

Participants will encounter interconnected investigation stages that may involve:
\begin{itemize}
    \item Reconnaissance and information gathering.
    \item Initial access and vulnerability identification.
    \item Enumeration and forensic investigation.
    \item Credential discovery and exploitation.
    \item Privilege escalation and lateral movement.
    \item Achievement of the final objective.
\end{itemize}

The challenge environments will be sandboxed and designed to provide a controlled
setting for practical cybersecurity experimentation.

\subsection*{Technical Resources and Learning Opportunities}

Alongside the competition, Veni Vidi Vici will feature \textbf{learning resources} aimed at
introducing participants to selected cybersecurity domains.

These resources will aim to:
\begin{itemize}
    \item Introduce fundamental concepts in different cybersecurity domains.
    \item Familiarize participants with commonly used security tools.
    \item Bridge the gap between theoretical understanding and hands-on cybersecurity
    practice.
\end{itemize}

The resources will be designed to accommodate participants with varying levels of
prior experience.

\section{Collaboration and Industry Involvement}

We invite organizations, cybersecurity companies, security researchers, and industry
professionals to collaborate with Veni Vidi Vici and contribute to the development
of practical cybersecurity education at IIT Dharwad.

Possible forms of collaboration include:
\begin{itemize}
    \item \textbf{Challenge Partnerships:} Contributing real-world security scenarios,
    vulnerabilities, or problem statements that can be adapted into competition
    challenges.
    \item \textbf{Technical Mentorship:} Providing cybersecurity professionals and
    researchers to mentor participants and guide technical sessions.
    \item \textbf{Knowledge Partnerships:} Collaborating on workshops, demonstrations,
    and discussions covering emerging areas of cybersecurity.
    \item \textbf{Judging and Evaluation:} Participating in the evaluation of challenges,
    technical presentations, and competition outcomes.
    \item \textbf{Infrastructure Partnerships:} Supporting the event through cybersecurity
    platforms, cloud infrastructure, software licenses, and technical resources.
    \item \textbf{Prize and Sponsorship Partnerships:} Supporting participant incentives,
    awards, and event execution.
\end{itemize}

Industry participation will provide organizations with an opportunity to engage with
technically motivated students, demonstrate their expertise, and contribute to
strengthening the cybersecurity ecosystem.

We also welcome opportunities to adapt industry-relevant security scenarios into
educational challenge environments, ensuring that participants gain exposure to
problems inspired by practical security research and professional workflows.

\section{Tentative Timeline}

\begin{center}
\small
\begin{tabularx}{\textwidth}{>{\RaggedRight\arraybackslash}p{0.23\textwidth}
                              >{\Centering\arraybackslash}p{0.22\textwidth}
                              >{\RaggedRight\arraybackslash}X}
\toprule
\textbf{Phase} & \textbf{Tentative Dates} & \textbf{Description}\\
\midrule
Event Announcement & \placeholder{19th December, 2026} &
Announcement of Veni Vidi Vici and opening of registrations.\\

Registration & \placeholder{20th December, 2026 -- 24th January, 2027} &
Participant registration and team formation.\\


Release of CTF Materials & \placeholder{20th January, 2027} &
Release of competition relevant materials, literature and resources.\\

Competition Launch & \placeholder{25th January, 2027} &
Commencement of the CTF competition and release of challenges.\\

Competition Conclusion & \placeholder{27th January, 2027} &
Closure of the competition and final evaluation.\\

Results and Awards & \placeholder{29th January, 2027} &
Announcement of results and recognition of winning participants.\\
\bottomrule
\end{tabularx}
\end{center}

\textit{The final schedule and competition format will be announced closer to the event.}

\section{Event Vision}

Veni Vidi Vici aspires to create more than a conventional cybersecurity competition.
It seeks to establish an environment where curiosity meets practical security research,
where theoretical knowledge is tested against realistic challenges, and where
participants develop the ability to understand and reason about complex security
problems.

By bringing together students, cybersecurity professionals, and industry organizations,
the event aims to promote awareness of emerging security threats, encourage technical
exploration, and nurture the next generation of cybersecurity practitioners.

We warmly welcome organizations and professionals working in cybersecurity,
information security, cloud security, application security, and related domains to
collaborate with us in making Veni Vidi Vici a meaningful technical experience.

\section*{Contact Details}

\begin{tabularx}{\textwidth}{@{}X X@{}}
\textbf{Ankush Tarafdar} & \textbf{PARSEC Organizing Team}\\
Event Lead -- Veni Vidi Vici & \textbf{Website:} \url{https://parsec.iitdh.ac.in}\\
Email: \href{mailto:mc25bt011@iitdh.ac.in}{mc25bt011@iitdh.ac.in} & Email: \href{mailto:events.parsec@iitdh.ac.in}{events.parsec@iitdh.ac.in}\\
\end{tabularx}

\vfill
\begin{center}
    {\color{VVVCyan}\rule{0.75\textwidth}{1pt}}\\[5pt]
    {\large\bfseries\color{VVVNavy} PARSEC 7.0 -- IIT DHARWAD}\\
    {\small Indian Institute of Technology Dharwad}
\end{center}

\end{document}